% 04_link.tex - the Raven files and the coupling file
\section{Link the Raven model to the MODFLOW model}

\step{Edit \texttt{Creek.rvh} and \texttt{Creek.rvi}}
These two changes are the same as for a groundwater model that Raven builds. In \cmd{Creek.rvh}, mark the coupled HRUs
with an aquifer profile and list them in the group \cmd{AquiferHRUs}:
\diff{t4_a_rvh.diff}
In \cmd{Creek.rvi}, switch the coupling on and send the percolation of the coupled HRUs to \cmd{GROUNDWATER} (Raven's store
for the water that becomes MODFLOW recharge):
\diff{t4_a_rvi.diff}

\step{Edit \texttt{Creek.rvp}}
The \cmd{.rvp} needs no aquifer classes: the layers and their properties come from the MODFLOW model. A profile is written
as \cmd{name, 0} (zero layers). It only marks the coupled HRUs and carries the exchange settings, here the conductance of
the \emph{seepage face} [1/d] (the land surface where the water table reaches it and groundwater seeps out):
\diff{t4_a_rvp.diff}

\step{Create \texttt{Creek.rvg}}\label{st:rvg}
Create the coupling file \cmd{Creek.rvg} with this content:
\lstinputlisting[style=file]{generated/t4_rvg_a.txt}
Or copy the file: \cmd{cp \$WS/model/step\_07/Creek.rvg .}
\why
\begin{itemize}
\item \cmd{:MF6Simulation} names the simulation file of the model. Raven copies the folder that holds it into
  \cmd{out/mf6/} and runs the copy. The folder \cmd{model} is never changed.
\item \cmd{:MF6CRS} gives the projection of the model's coordinates. Raven reads the cell outlines from the model
  (\cmd{DELR}, \cmd{DELC}, \cmd{XORIGIN}, \cmd{YORIGIN}, \cmd{ANGROT}), converts feet to metres, and places them on the
  map under the HRU polygons. For a model without a projection, \cmd{:MF6GridFile} (a file of cell polygons) or
  \cmd{:OverlapWeights} (a table of the overlaps of HRUs and cells) link them instead.
\item \cmd{:MF6StartDate} is the calendar date of the model's time zero. Raven keeps the model's stress periods and splits
  each into Raven days.
\item \cmd{:RavenTakesOver RCH\_MODEL EVT\_MODEL}: every recharge, ET or UZF package of the model must be named, either as
  taken over (removed; Raven supplies the recharge) or as kept (\cmd{:KeepPackages}). A package that is not named stops
  the run, so the recharge is never counted twice.
\item \cmd{:StreamPackage RIV\_CREEK AUX SUBBASIN}: the river package whose flows go to Raven's river reaches, and the
  auxiliary variable that gives the subbasin of each cell.
\item \cmd{:SeepageToRaven ON} adds a drain (\cmd{DRN\_RAVEN}) at the top of the coupled columns. The water it drains goes
  back to Raven.
\end{itemize}
